\documentclass[12pt]{article}
\usepackage[utf8]{inputenc}
\usepackage{graphicx}
\usepackage{amsmath}
\usepackage{amssymb}
\usepackage{mathrsfs}
\usepackage{pgfornament}
\usepackage[many]{tcolorbox}
\usepackage[margin = 0.5in]{geometry}
\usepackage{collectbox}
\usepackage{mdframed}
\usepackage{xcolor}
\usepackage{mathtools}
\DeclarePairedDelimiter\ceil{\lceil}{\rceil}
\DeclarePairedDelimiter\floor{\lfloor}{\rfloor}
\newcommand\norm[1]{\left\lVert#1\right\rVert}
\newcommand\textbox[1]{%
  \parbox{.333\textwidth}{#1}%
}

\linespread{2}

\makeatletter
\newcommand{\mybox}{%
    \collectbox{%
        \setlength{\fboxsep}{2pt}%
        \fbox{\BOXCONTENT}%
    }%
}
\makeatother

\usepackage{listings}
\usepackage{color}

\definecolor{dkgreen}{rgb}{0,0.6,0}
\definecolor{gray}{rgb}{0.5,0.5,0.5}
\definecolor{mauve}{rgb}{0.58,0,0.82}

\lstset{frame=tb,
  language=R,
  aboveskip=3mm,
  belowskip=3mm,
  showstringspaces=false,
  columns=flexible,
  basicstyle={\small\ttfamily},
  numbers=none,
  numberstyle=\tiny\color{gray},
  keywordstyle=\color{blue},
  commentstyle=\color{dkgreen},
  stringstyle=\color{mauve},
  breaklines=true,
  breakatwhitespace=true,
  tabsize=3
}
\title{MAT 524 HW 3}
\begin{document}

\hfill Feb. 20, 2019
\begin{center}
\begin{large}
\textbf{MAT 524: Functions of a Real Variable II \\
\vspace{-5mm}
Homework III\\}
\end{large}
\vspace{-2mm}
Nolan R. H. Gagnon \\
\end{center}

\vspace{5mm}

\noindent \textbf{{\huge [}{\Large [}} Problem No. 1 \textbf{{\Large ]}{\huge ]}}{\Large \textbf{:}}

\vspace{3mm}

\noindent\mybox{\colorbox{yellow}{\textbf{Problem Statement}}}\hspace{3mm}\hrulefill

\noindent Suppose $f_n, f\in L^1(E)$. We say that $f_n\rightarrow f$ in $L^1$ if $\norm{f-f_n}_1\rightarrow 0$, i.e., $$\int_E|f_n-f| \longrightarrow 0.$$
Compare the two conditions

\noindent(a) $f_n\rightarrow f$ almost everywhere\\
\noindent(b) $f_n\rightarrow f$ in $L^1.$

\noindent Determine whether (a) implies (b) or (b) implies (a). Give a proof or counter-example for each implication.

\noindent\mybox{\colorbox{yellow}{\textbf{Solution}}}\hspace{3mm}\hrulefill

\noindent \textbf{Counter-example for (a) $\Rightarrow$ (b):} Suppose $E=[0,1]$ and $f_n(x) = nx^2(1-x^3)^n$. For $x\in E$, $$\lim\limits_{n\rightarrow\infty} nx^2(1-x^3)^n = 0$$ because $a^n$ shrinks faster than $n$ grows for $0\leq a < 1$. So $f_n\rightarrow 0$ almost everywhere, but
\begin{align}
    \lim\limits_{n\rightarrow\infty}\int_E \left|nx^2(1-x^3)^n - 0\right| \hspace{1mm} dx &= \lim\limits_{n\rightarrow\infty}\frac{-n}{3}\int_E -3x^2(1-x^3)^n \hspace{1mm} dx \\
    &= \lim\limits_{n\rightarrow\infty} \left[\frac{-n(1-x^3)^{n+1}}{3n+3}\right]_0^1 \\
    &= \frac{1}{3}.
\end{align}
Thus, $f_n$ doesn't converge to $f$ in $L^1$.

\vspace{3mm}

\noindent\textbf{Counter-example for (b) $\Rightarrow$ (a):} Suppose $E=[0,1]$ and $f_n = \sin(nx)$. Note that$f_n$ converges to $0$ in $L^1$ because 
\begin{align}
    \lim\limits_{n\rightarrow\infty}\int_E |\sin(nx)-0| \hspace{1mm} dx &= \lim\limits_{n\rightarrow\infty}\int_E \sin(nx) \hspace{1mm} dx \\
    &= \lim\limits_{n\rightarrow\infty}\left[\frac{-\cos(n)}{n}+\frac{1}{n}\right] \\
    &= 0.
\end{align}
But $f_n=\sin(nx)$ does not converge a.e. on $E$ as $n\rightarrow\infty$, due to the oscillating nature of the sine function.

\hfill$\blacksquare$

\vfill

\begin{center}
\pgfornament[width = 50mm, color = black, symmetry = h]{83}\\
\vspace{2mm}
\vspace{-0.65cm}
\hrulefill \\
\vspace{-0.95cm}
\hrulefill
\end{center}

\pagebreak

\noindent \textbf{{\huge [}{\Large [}} Problem No. 2 \textbf{{\Large ]}{\huge ]}}{\Large \textbf{:}}

\vspace{3mm}

\noindent\mybox{\colorbox{yellow}{\textbf{Problem Statement}}}\hspace{3mm}\hrulefill 

\noindent Let $f_n,f$ be integrable measurable functions on $E$. Assume $f_n\rightarrow f$ almost everywhere. Suppose that $\lim\limits_{n\rightarrow\infty}\int_E |f_n| = \int_E |f|$. Prove that $f_n\rightarrow f$ in $L^1(E)$.

\vspace{5mm}

\noindent\mybox{\colorbox{yellow}{\textbf{Solution}}}\hspace{3mm}\hrulefill

\noindent By the Triangle Inequality, $|f_n - f| \leq |f_n| + |f|$ for all $n\in\mathbb{N}$. Furthermore, since $\lim\limits_{n\rightarrow\infty}\int_E |f_n| = \int_E |f|$ and $f$ is integrable, we have 
\begin{align}
    \lim\limits_{n\rightarrow\infty}\int_E |f_n| + \int_E |f| &= \int_E |f| + \int_E |f|,
\end{align}
i.e.,
\begin{align}
    \lim\limits_{n\rightarrow\infty}\int_E \left(|f_n| + |f|\right) &= 2\int_E |f| \\
    & < \infty.
\end{align}
Thus, by the Generalized Dominated Convergence Theorem (GDCT), 
\begin{align}
    \lim\limits_{n\rightarrow\infty}\int_{E} |f_n - f| &= \int_{E} \lim\limits_{n\rightarrow\infty} |f_n-f| \\
    &= 0.
\end{align}
Therefore, $f_n$ converges to $f$ in $L^1(E)$.

\hfill$\blacksquare$

\vfill

\begin{center}
\pgfornament[width = 50mm, color = black, symmetry = h]{83}\\
\vspace{2mm}
\vspace{-0.65cm}
\hrulefill \\
\vspace{-0.95cm}
\hrulefill
\end{center}

\pagebreak

\noindent \textbf{{\huge [}{\Large [}} Problem No. 3 \textbf{{\Large ]}{\huge ]}}{\Large \textbf{:}}

\vspace{3mm}

\noindent\mybox{\colorbox{yellow}{\textbf{Problem Statement}}}\hspace{3mm}\hrulefill 

\noindent Let $C([0,1])$ denote the space of continuous real-valued functions on the interval $[0,1]$. Show that it is complete relative to the max-norm.

\vspace{5mm}

\noindent\mybox{\colorbox{yellow}{\textbf{Solution}}}\hspace{3mm}\hrulefill

\noindent Let $\lbrace f_n \rbrace$ be a Cauchy sequence of continuous functions in $C([0,1])$. Since $C([0,1])$ is a subspace of $L^{\infty}([0,1])$ and $[0,1]$ is measurable, we can use Proposition 5 and Theorem 7 to conclude that $\lbrace f_n\rbrace$ has a rapidly Cauchy subsequence that converges. Thus, by Proposition 4, $\lbrace f_n \rbrace$ converges. Let us call the function to which $\lbrace f_n \rbrace$ converges $f$. To show that $C([0,1])$ is complete with respect to the max-norm, we must show that $f\in C([0,1])$. In order to achieve this, we will show that $\norm{f_n-f}_{\infty} \rightarrow 0$. 

\hfill$\blacksquare$

\vfill

\begin{center}
\pgfornament[width = 50mm, color = black, symmetry = h]{83}\\
\vspace{2mm}
\vspace{-0.65cm}
\hrulefill \\
\vspace{-0.95cm}
\hrulefill
\end{center}

\pagebreak

\noindent \textbf{{\huge [}{\Large [}} Problem No. 4 \textbf{{\Large ]}{\huge ]}}{\Large \textbf{:}}

\vspace{3mm}

\noindent\mybox{\colorbox{yellow}{\textbf{Problem Statement}}}\hspace{3mm}\hrulefill 

\noindent Let $f$ be a bounded measurable function on $[0,1]$. Prove that $$\lim\limits_{p\rightarrow\infty} \norm{f}_p = \norm{f}_{\infty}.$$

\vspace{5mm}

\noindent\mybox{\colorbox{yellow}{\textbf{Solution}}}\hspace{3mm}\hrulefill

\noindent 

\hfill$\blacksquare$

\vfill

\begin{center}
\pgfornament[width = 50mm, color = black, symmetry = h]{83}\\
\vspace{2mm}
\vspace{-0.65cm}
\hrulefill \\
\vspace{-0.95cm}
\hrulefill
\end{center}

\pagebreak

\noindent \textbf{{\huge [}{\Large [}} Problem No. 5 \textbf{{\Large ]}{\huge ]}}{\Large \textbf{:}}

\vspace{3mm}

\noindent\mybox{\colorbox{yellow}{\textbf{Problem Statement}}}\hspace{3mm}\hrulefill 

\noindent Let $1\leq p < \infty$. Suppose $f_n, f\in L^p(E)$ and $f_n\rightarrow f$ pointwise almost everywhere. Show that $\norm{f_n}_p\rightarrow \norm{f}_p$ if and only if $f_n\rightarrow f$ in $L^p(E)$.

\vspace{5mm}

\noindent\mybox{\colorbox{yellow}{\textbf{Solution}}}\hspace{3mm}\hrulefill

\noindent For the reverse direction of the proof, we need to use the $p$-norm analogue of the reverse triangle inequality. We have the following lemma:

\vspace{5mm}

\noindent\textbf{Lemma 5.1:} If $f,g$ are in $L^p(E)$, then $\left|\norm{f}_p-\norm{g}_p\right| \leq \norm{f-g}_p.$

\vspace{5mm}

\noindent\textbf{Proof:} By Minkowski's Inequality, we have $$\norm{f+g}_p \leq \norm{f}_p + \norm{g}_p.$$ Using this, we obtain the following inequalities: 
\begin{align}
    \norm{f}_p &= \norm{f-g+g}_p \leq \norm{f-g}_p + \norm{g}_p \label{mink1}\\ 
    \norm{g}_p &= \norm{g-f+f}_p \leq \norm{g-f}_p + \norm{f}_p \label{mink2}
\end{align}
Rearranging (\ref{mink1}) and (\ref{mink2}), we see that
\begin{align}
    \norm{f}_p - \norm{g}_p &\leq \norm{f-g}_p \\
    \norm{g}_p - \norm{f}_p &\leq \norm{g-f}_p.
\end{align}
But $\norm{f-g}_p=\norm{g-f}_p$, so we conclude that $$\left|\norm{f}_p-\norm{g}_p\right| \leq \norm{f-g}_p.$$

\hfill$\blacksquare$

\noindent Now we can move on the the proof of the result in the problem statement.

\pagebreak

\noindent\textbf{Proof of Main Result:} \textcolor{blue}{\textbf{($\Rightarrow$)}} First suppose that $\norm{f_n}_p\rightarrow \norm{f}_p$. Then 

\begin{align}
\lim\limits_{n\rightarrow\infty}\left[\int_E |f_n|^p\right]^{\frac{1}{p}} &= \left[\int_E|f|^p\right]^{\frac{1}{p}} \label{eq5.1}
\end{align}
Taking the $p$th power of both sides of Equation (\ref{eq5.1}) and applying the limit definition of continuity, we see that
\begin{align}
    \lim\limits_{n\rightarrow\infty}\int_E|f_n|^p &= \int_E|f|^p. \label{eq5.2}
\end{align}
Adding $\int_E |f|^p$ (which is finite, since $f\in L^p(E)$) to both sides of Equation (\ref{eq5.2}) and multiplying the result by $2^p$, we obtain
\begin{align}
    \lim\limits_{n\rightarrow\infty} \int_E 2^p\left(|f_n|^p+|f|^p\right) &= 2^{p+1}\int_E |f|^p \\
    &< \infty.
\end{align}
We also have that $|f_n-f|^p \leq 2^p\left(|f_n|^p+|f|^p\right).$ Thus, by the GDCT, \begin{align}
    \lim\limits_{n\rightarrow\infty}\left[\int_E|f_n-f|^p\right]^\frac{1}{p} &= \left[\lim\limits_{n\rightarrow\infty}\int_E|f_n-f|^p\right]^\frac{1}{p} \\
    &= \left[\int_E\lim\limits_{n\rightarrow\infty}|f_n-f|^p\right]^\frac{1}{p} \\
    &= 0^{\frac{1}{p}} \\
    &= 0.
\end{align}
Therefore, $f_n\rightarrow f$ in $L^p(E).$

\vspace{5mm}

\noindent \textcolor{blue}{\textbf{($\Leftarrow$)}} Next, suppose that $f_n\rightarrow f$ in $L^p(E)$. By Lemma 5.1, 
\begin{align}
    \left|\norm{f_n}_p-\norm{f}_p\right| \leq \norm{f_n-f}_p. \label{eq5.3}
\end{align}
\noindent Taking the limit on both sides of (\ref{eq5.3}) yields
\begin{align}
    \lim\limits_{n\rightarrow\infty}\left|\norm{f_n}_p-\norm{f}_p\right| \leq 0 \label{eq5.4},
\end{align}
since $f_n\rightarrow f$ in $L^p(E)$. But the right-hand side of (\ref{eq5.4}) must be non-negative, so it follows that 
\begin{align}
    \lim\limits_{n\rightarrow\infty}\left|\norm{f_n}_p-\norm{f}_p\right| &= 0,
\end{align}
i.e.,
\begin{align}
    \lim\limits_{n\rightarrow\infty}\norm{f_n}_p &= \norm{f}_p,
\end{align}
as desired.

\hfill$\blacksquare$ 

\vfill

\begin{center}
\pgfornament[width = 50mm, color = black, symmetry = h]{83}\\
\vspace{2mm}
\vspace{-0.65cm}
\hrulefill \\
\vspace{-0.95cm}
\hrulefill
\end{center}

\end{document}
